\section{Discussion}
\label{sec:Discussion}

For the calculation of the relative experimental error, the formula
\begin{equation}
    \Delta x =   \frac{ \left| x_{exp} - x_{theo} \right|   }{x_{theo}} \, ,
    \label{eq:abweichung}
\end{equation}
is being used.

\subsection{Detector calibration with \texorpdfstring{\ce{^{152}Eu}}{152Eu}}

The linear regression \autoref{eqn:lin} gave the parameters
\begin{align}
    m &= \SI{0.11950 \pm 0.00005}{keV}\text{/channel} \\
    b &= \SI{-0.98 \pm 0.23}{keV} \ .
\end{align}

$b$ makes physically no sense, because it would mean that at an energy of 0 the detector emits energy, but it is so small that it can be neglected.
To check the calibration, the energies of the \ce{^{152}Eu} lines calculated with \autoref{eqn:lin} from the fitted peak positions are compared with the literature values in \autoref{tab:disc_eu}.
In this it is visible that the error is smaller than $\SI{0.1}{\%}$, so the energy calibration fits very well for those values and can be used for all the other steps.
\begin{table}[H]
	\centering
	\caption{Literature energies $E_\text{lit}$~\cite{eu152} of the \ce{^{152}Eu} lines compared with the energies $E_\text{exp}$ from the gaussian fits and the calibration.}
	\label{tab:disc_eu}
	\begin{tabular}{c  c  c}
		\toprule
		$E_\text{lit}$ / keV & $E_\text{exp}$ / keV & $\Delta E$ / \%\\
		\midrule
		121.78 & 121.80  \, \clap{$\pm$} \,  0.24 & 0.015 \\
		 244.70 & 244.70  \, \clap{$\pm$} \,  0.25 & 0.001 \\
		 344.28 & 344.34  \, \clap{$\pm$} \,  0.27 & 0.018 \\
		 411.12 & 411.11  \, \clap{$\pm$} \,  0.29 & 0.003 \\
		 443.96 & 444.01  \, \clap{$\pm$} \,  0.30 & 0.010 \\
		 778.90 & 779.01  \, \clap{$\pm$} \,  0.41 & 0.014 \\
		 867.37 & 867.48  \, \clap{$\pm$} \,  0.45 & 0.012 \\
		 \bottomrule
	\end{tabular}
\end{table}


\subsection{Monochromatic gamma source \texorpdfstring{\ce{^{137}Cs}}{137Cs}}


The measured positions of the photopeak, the backscatter peak and the compton edge are compared with the literature and theory values
\begin{align}
    E_\text{photo, exp} &= \SI{661.6 \pm 0.4}{keV} & E_\text{photo, lit} &= \SI{661.657}{keV} & \Delta E_\text{photo} &= \SI{0.007}{\percent} \\
    E_\text{back, exp} &= \SI{188.9 \pm 0.4}{keV} & E_\text{back, theo} &= \SI{184.3}{keV} & \Delta E_\text{back} &= \SI{2.5}{\percent} \\
    E_\text{ce, exp} &= \SI{477.0 \pm 0.4}{keV} & E_\text{ce, theo} &= \SI{477.3}{keV} & \Delta E_\text{ce} &= \SI{0.06}{\percent} \ .
\end{align}
The photopeak fits very well to the literature value, so the calibration also works for this source.
The full width at half maximum of $\SI{2.386 \pm 0.020}{keV}$ shows the very good energy resolution of the germanium detector.
\\
The backscatter peak is $\SI{2.5}{\percent}$ above the theory value.
This can be explained, because the peak is very broad and lies on the compton continuum, so the gaussian fit describes it only roughly.
In addition the gamma rays are not only scattered at exactly $\SI{180}{\degree}$ but in a range of angles around it, which shifts the peak to higher energies.
\\
The compton edge was determined by eye, because it is smeared out by the resolution of the detector and by multiple compton scattering, so it cannot be fitted well.
For this an error of $\pm 3$ channels is assumed, which is $\pm\SI{0.4}{keV}$ with the calibration.
Within this the compton edge fits to the theory value.

The ratio of the line contents of the photopeak and the compton continuum is compared with the ratio of the absorption probabilities
\begin{align}
    \frac{N_\text{photo}}{N_\text{compton}} &= \num{0.2784 \pm 0.0030} & \frac{P_\text{photo}}{P_\text{compton}} &= \num{0.040} & \Delta &= \SI{592}{\percent} \ .
\end{align}
The measured ratio is about 7 times bigger than the ratio of the absorption probabilities.
The reason is that $P$ only describes the first interaction of the gamma ray in the detector.
A lot of gamma rays are first compton scattered and then fully absorbed in the crystal, so they end up in the photopeak and not in the compton continuum.
Also the extinction coefficients are only read from a plot and the compton fit is only done from $\SI{350}{keV}$ to $\SI{460}{keV}$ and then extrapolated to the whole continuum,
which both make the result more inaccurate.

\subsection{Activity determination of \texorpdfstring{\ce{^{133}Ba}}{133Ba}}

The measured energies of the \ce{^{133}Ba} lines are compared with the literature values in \autoref{tab:disc_ba}.


\begin{table}[H]
	\centering
	\caption{Literature energies $E_\text{lit}$~\cite{ba133} of the \ce{^{133}Ba} lines compared with the measured energies $E_\text{exp}$.}
	\label{tab:disc_ba}
	\begin{tabular}{c  c  c}
		\toprule
		$E_\text{lit}$ / keV & $E_\text{exp}$ / keV & $\Delta E$ / \%\\
		\midrule
		276.40 & 276.41  \, \clap{$\pm$} \,  0.26 & 0.005 \\
		 302.85 & 302.89  \, \clap{$\pm$} \,  0.26 & 0.012 \\
		 356.02 & 356.05  \, \clap{$\pm$} \,  0.28 & 0.007 \\
		 383.85 & 383.85  \, \clap{$\pm$} \,  0.28 & 0.000 \\
		 \bottomrule
	\end{tabular}
\end{table}


Once again the energy values are very close to the literature values, so it is shown again that the energy calibration fits.
\\
The calculated mean activity of the \ce{^{133}Ba} source is
\begin{align}
    A_{\ce{^{133}Ba}} = \SI{771 \pm 10}{Bq} \ .
\end{align}
The activities of the single lines are between $\SI{750}{Bq}$ and $\SI{793}{Bq}$ and differ more from each other than their errors.
This could be because of a systematic error in the fit of $Q$.

\subsection{Gamma spectrum of an unknown source}

The measured energies of the unknown source are compared with the literature values of the assigned nuclides in \autoref{tab:disc_unknown}.

\begin{table}[H]
	\centering
	\caption{Literature energies $E_\text{lit}$~\cite{nndc} of the assigned nuclides compared with the measured energies $E_\text{exp}$ of the unknown source.}
	\label{tab:disc_unknown}
	\begin{tabular}{c  c  c  c}
		\toprule
		Nuclide & $E_\text{lit}$ / keV & $E_\text{exp}$ / keV & $\Delta E$ / \%\\
		\midrule
		\ce{^{226}Ra} & 186.21 & 185.96  \, \clap{$\pm$} \,  0.24 & 0.134 \\
		 \ce{^{214}Pb} & 242.00 & 242.00  \, \clap{$\pm$} \,  0.25 & 0.000 \\
		 \ce{^{214}Pb} & 295.22 & 295.21  \, \clap{$\pm$} \,  0.26 & 0.006 \\
		 \ce{^{214}Pb} & 351.93 & 351.91  \, \clap{$\pm$} \,  0.28 & 0.007 \\
		 \ce{^{214}Bi} & 609.32 & 609.31  \, \clap{$\pm$} \,  0.35 & 0.001 \\
		 \bottomrule
	\end{tabular}
\end{table}


The calculated activities of the unknown source are
\begin{align}
    A_{\ce{^{226}Ra}} &= \SI{1750 \pm 70}{Bq} \\
    A_{\ce{^{214}Pb}} &= \SI{1216 \pm 41}{Bq} \\
    A_{\ce{^{214}Bi}} &= \SI{1400 \pm 40}{Bq} \ .
\end{align}
Because \ce{^{214}Pb} and \ce{^{214}Bi} are daughter nuclides of \ce{^{226}Ra}, they should all have the same activity in the secular equilibrium.
\ce{^{214}Pb} and \ce{^{214}Bi} differ by about $\SI{14}{\percent}$, which can be explained by the uncertainty of $Q(E)$ and of the distance $a$.
Also the activities of the three \ce{^{214}Pb} lines increase with the energy, which also shows that $Q(E)$ is not exact for these energies.
This all leads to the conclusion that it is a uranium source.

In general it can be said that the germanium detector is very well suited to determine gamma spectra of different sources.

\clearpage
